\documentclass[11pt]{article}
\usepackage[T1]{fontenc}
\usepackage{lmodern,amsmath,amssymb,amsthm}
\usepackage[margin=1in]{geometry}
\usepackage{hyperref}

\newtheorem{theorem}{Theorem}
\newtheorem{lemma}[theorem]{Lemma}
\newtheorem{corollary}[theorem]{Corollary}

\title{Effective Detection of Mordell--Weil Classes
by Covector-Line Kummer Witnesses}
\author{\small \href{mailto:admin@cosmosgeek.org}{admin@cosmosgeek.org}\qquad \href{mailto:b321014@iiit-bh.ac.in}{b321014@iiit-bh.ac.in}}
\date{Working paper}

\begin{document}
\maketitle

\begin{abstract}
Under an explicit full first-level Kummer-image
hypothesis over the infinite torsion field, $r$
distinct good rational primes of size at most
$\exp(C\ell^n\log(2\ell^n))$ suffice to detect every
nonzero class of a subgroup generated by $r$ points
modulo $\ell^n$.

The proof uses a primitive covector whose projective
stabilizer has index $O_\ell(\ell^n)$, additive Kummer
characters over the corresponding number field,
and Thorner--Zaman's unconditional degree-one
Chebotarev bound.

Using one coefficient vector at every prime gives an averaged reduction formula for the rank of the specified subgroup.
\end{abstract}

\section{Statement}

Let $E/\mathbb Q$ be an elliptic curve, and let
\[
\Gamma=\langle P_1,\ldots,P_r\rangle
\subset E(\mathbb Q),
\qquad r\ge1,
\]
where the $P_i$ are $\mathbb Z$-linearly independent.

Fix a rational prime $\ell$, and write
\[
K_\infty=\mathbb Q(E[\ell^\infty]).
\]

We impose the following explicit hypothesis.

\medskip
\noindent
\textbf{Full first-level Kummer hypothesis
$(\mathrm K_\ell)$.}
Choose $R_i\in E(\overline{\mathbb Q})$ with
$[\ell]R_i=P_i$. The map
\[
\begin{aligned}
\kappa_\ell:
\operatorname{Gal}(\overline{\mathbb Q}/K_\infty)
&\longrightarrow E[\ell]^r,\\
\sigma&\longmapsto
(\sigma R_i-R_i)_{i=1}^r
\end{aligned}
\]
is surjective.

Because $K_\infty$ contains $E[\ell]$, this is a
homomorphism. Its image is independent of the
chosen division points.

This hypothesis is stronger than independence of
the $P_i$ in $\Gamma/\ell\Gamma$. It is also stronger
than merely knowing full first-level Kummer image
over $\mathbb Q(E[\ell])$: intersection with higher
torsion fields must be controlled.

For a good rational prime $p\ne\ell$ and
$m=\ell^n$, define
\[
\Phi_{m,p}:
(\mathbb Z/m\mathbb Z)^r
\longrightarrow E(\mathbb F_p)/mE(\mathbb F_p)
\]
by
\[
\Phi_{m,p}(a_1,\ldots,a_r)
=\sum_{i=1}^r a_i\overline P_i
\bmod mE(\mathbb F_p).
\]

For a finite set $\mathcal P$ of good primes,
write $\Phi_{m,\mathcal P}$ for the product
of these maps.

\begin{theorem}[Effective covector-line detection]
\label{thm:main}
Assume $(\mathrm K_\ell)$.

There exists a constant
$C=C(E,\Gamma,\ell)>0$ such that, for every
$n\ge1$ and $m=\ell^n$, there are $r$
distinct rational primes
\[
p_1,\ldots,p_r
\le \exp\!\bigl(Cm\log(2m)\bigr),
\]
all of good reduction for $E$ and different
from $\ell$, for which
\[
\boxed{
\Phi_{m,\{p_1,\ldots,p_r\}}
\text{ is injective}.
}
\]

In particular, every nonzero
$a\in(\mathbb Z/m\mathbb Z)^r$ satisfies
\[
\sum_i a_i\overline P_i
\notin mE(\mathbb F_{p_j})
\]
for at least one $j$.

The result holds without GRH and without an
assumption that the full linear image
$\rho_{E,m}(G_{\mathbb Q})$ equals
$\operatorname{GL}_2(\mathbb Z/m\mathbb Z)$.
\end{theorem}

\section{The covector field}

Put
\[
R=\mathbb Z/m\mathbb Z,\qquad
T=E[m]\simeq R^2,
\]
and let
\[
H_m=\rho_{E,m}(G_{\mathbb Q})
\subseteq\operatorname{GL}_2(R).
\]

Choose a primitive surjective covector
\[
\lambda:T\longrightarrow R.
\]

Thus the reduction of $\lambda$ modulo
$\ell$ is nonzero.

Define
\[
H_\lambda=
\left\{
M\in H_m:
\lambda M=u(M)\lambda,\
u(M)\equiv1\pmod\ell
\right\}.
\]
The scalar $u(M)\in R^\times$ is unique,
and
\[
u:H_\lambda\longrightarrow R^\times
\]
is a group homomorphism.

Let
\[
F=\mathbb Q(E[m])^{H_\lambda}.
\]
This field need not be Galois over
$\mathbb Q$.

\begin{lemma}
\label{lem:degree}
The degree $d=[F:\mathbb Q]$ satisfies
\[
\boxed{
d\le(\ell^2-1)\ell^{n-1}
=O_\ell(m).
}
\]
\end{lemma}

\begin{proof}
The primitive rank-one direct summands of
$R^2$ are parametrized by
$\mathbb P^1(R)$, whose cardinality is
\[
|\mathbb P^1(R)|=(\ell+1)\ell^{n-1}.
\]

The subgroup of $\operatorname{GL}_2(R)$
preserving $R\lambda$ therefore has index
$(\ell+1)\ell^{n-1}$.

Within that stabilizer, requiring the
scalar action on $\lambda$ to be
$1$ modulo $\ell$ increases the index
by at most $\ell-1$.

Intersecting this subgroup with $H_m$
cannot increase its index relative to
$H_m$. Hence
\[
[H_m:H_\lambda]
\le(\ell+1)(\ell-1)\ell^{n-1}.
\]
This equals $[F:\mathbb Q]$.
\end{proof}

The reduction of the scalar character
$u$ is trivial, although $u$ itself
need not be trivial.

This distinction is essential.

\section{Additive Kummer characters}

Choose $Q_i\in E(\overline{\mathbb Q})$
satisfying
\[
[m]Q_i=P_i
\]
and put
\[
\kappa_{m,i}(\sigma)
=\sigma Q_i-Q_i\in T.
\]

The Kummer cocycle relation is
\[
\kappa_{m,i}(\sigma\tau)
=
\kappa_{m,i}(\sigma)
+\rho_m(\sigma)\kappa_{m,i}(\tau).
\]

For $\sigma\in G_F$, define
\[
\psi_i(\sigma)
=
\lambda(\kappa_{m,i}(\sigma))
\bmod\ell.
\]

\begin{lemma}
\label{lem:characters}
Each
\[
\psi_i:G_F\longrightarrow\mathbb F_\ell
\]
is a continuous additive character.

Moreover, the combined map
\[
(\psi_1,\ldots,\psi_r):
\ker(u\circ\rho_m|_{G_F})
\longrightarrow\mathbb F_\ell^r
\]
is surjective.
\end{lemma}

\begin{proof}
For $\sigma\in G_F$,
\[
\lambda\rho_m(\sigma)
=u(\sigma)\lambda
\equiv\lambda\pmod\ell.
\]

Applying $\lambda$ modulo $\ell$ to
the cocycle relation gives
\[
\psi_i(\sigma\tau)
=\psi_i(\sigma)+\psi_i(\tau).
\]

Thus $\psi_i$ is a homomorphism.
This is not generally true for the
unreduced $R$-valued function
$\lambda\kappa_{m,i}$.

Multiplication by $\ell^{n-1}$ induces
a canonical isomorphism
\[
T/\ell T\simeq E[\ell].
\]
Under it, the reduction of
$\kappa_{m,i}$ corresponds to the
first-level Kummer cocycle for $P_i$.

The restriction to
$G_{K_\infty}$ has trivial linear
action and hence trivial $u$.
Hypothesis $(\mathrm K_\ell)$ says that
the joint first-level translations
fill $E[\ell]^r$.

Applying the nonzero functional induced
by $\lambda$ gives all of
$\mathbb F_\ell^r$. Therefore the same
surjectivity holds on the larger
kernel of $u$.
\end{proof}

The lemma exhibits a finite abelian
quotient of $G_F$ through
\[
\sigma\longmapsto
\bigl(u(\sigma),\psi_1(\sigma),
\ldots,\psi_r(\sigma)\bigr).
\]

Its image contains every pattern
\[
\boxed{
u=1,\qquad
(\psi_1,\ldots,\psi_r)=e_i,
\quad 1\le i\le r,
}
\]
where $e_i$ is the $i$th standard
basis vector.

\section{The local detection criterion}

\begin{lemma}
\label{lem:detection}
Suppose $p$ is a good rational prime,
$p\ne\ell$, unramified in the division
extension. Suppose a degree-one prime
of $F$ above $p$ has Frobenius satisfying
\[
u(\operatorname{Frob})=1,
\qquad
(\psi_1,\ldots,\psi_r)
(\operatorname{Frob})=e_i.
\]

Then every coefficient vector
$a=(a_1,\ldots,a_r)\in R^r$ for which
$a_i$ has minimum $\ell$-adic valuation
among the nonzero coordinates is detected
at $p$:
\[
\Phi_{m,p}(a)\ne0.
\]
\end{lemma}

\begin{proof}
Let $a\ne0$. Write
\[
a=\ell^k b,\qquad 0\le k<n,
\]
where at least one $b_j$ is a unit
modulo $\ell$.

For the index $i$ specified in the lemma, we have $b_i\not\equiv0\pmod\ell$. Choose a prime of $\overline{\mathbb Q}$ above the stated degree-one prime of $F$. Since $p$ is also unramified in the normal closure of $F$, its decomposition group lies in the corresponding conjugate of $G_F$; we evaluate Frobenius with that choice.

Write
\[
M=\rho_m(\operatorname{Frob}_p)
\]
and let
\[
\kappa_a=\sum_j a_j\kappa_{m,j}
(\operatorname{Frob}_p).
\]

Since $u(M)=1$,
\[
\lambda(M-I)=0\quad\text{in }R.
\]

On the other hand, the prescribed
character pattern gives
\[
\begin{aligned}
\lambda(\kappa_a)
&=\ell^k
\sum_j b_j\lambda(\kappa_{m,j})\\
&\equiv\ell^k b_i
\not\equiv0\pmod{\ell^{k+1}}.
\end{aligned}
\]

In particular,
\[
\kappa_a\notin(M-I)T.
\]

The reduction of
$P(a)=\sum_j a_jP_j$ has an
$\mathbb F_p$-rational $m$-division
point if and only if
\[
\kappa_a\in(M-I)T.
\]
Indeed, changing a geometric division
point by $T\in E[m]$ changes its
Frobenius displacement by $(M-I)T$.

The displayed noncontainment therefore
proves
\[
P(a)_p\notin mE(\mathbb F_p).
\]
\end{proof}

Notice that only the character pattern
modulo $\ell$ was needed. The condition
$u=1$ must hold modulo the entire
$m=\ell^n$.

\section{Avoiding bad primes and obtaining
distinct witnesses}

Let $S_0$ contain the rational primes
dividing $\ell N_E$, where $N_E$ is the
conductor of $E$.

Choose once and for all an auxiliary
prime $q\notin S_0$ sufficiently large
that there exist $r$ distinct nonzero
residue classes
\[
c_1,\ldots,c_r\pmod q
\]
none occupied by a prime in $S_0$.

For example, it suffices to take $q$
larger than every prime in $S_0$ and
with
\[
q-1>r+|S_0|.
\]

The torsion and Kummer extensions are
unramified at $q$: $E$ has good reduction
there and $m$ is prime to $q$.

Consequently the abelian extension of
$F$ cut out by $u$ and the $\psi_i$
is unramified above $q$.

The field $F$ is also unramified at $q$.
Thus $F(\zeta_q)/F$ is totally ramified
at every prime above $q$, and every
nontrivial subextension is ramified
there.

It follows that $F(\zeta_q)$ is linearly
disjoint, over $F$, from the abelian
Kummer-character extension.

Hence the simultaneous conditions
\[
\boxed{
u=1,\quad
(\psi_1,\ldots,\psi_r)=e_i,\quad
p\equiv c_i\pmod q
}
\]
define nonempty classes in one finite
abelian extension $L/F$.

The prescribed residue classes ensure
that their rational prime witnesses
are distinct and avoid $S_0$.

The full construction is abelian
over $F$, even though $F/\mathbb Q$
need not be Galois.

\section{Uniform discriminant and conductor bounds}

We now bound the quantities entering
the unconditional least-prime theorem.

Let
\[
d=[F:\mathbb Q]\ll_\ell m.
\]

The field $F$ is a subfield of
$\mathbb Q(E[m])$, so it is unramified
outside the fixed set $S_0$.

For a rational prime $p\in S_0$ and
$v\mid p$ in $F$, let $e_v,f_v$
denote its ramification and residue
degrees.

The standard local different bound,
together with $e_v\le d$, yields
\[
v_p(D_F)\ll d(1+\log_p d).
\]

Since $S_0$ is fixed,
\[
\boxed{
\log|D_F|\ll_{E,\ell}d\log(2d)
\ll_{E,\ell}m\log(2m).
}
\tag{1}
\]

Let
\[
Q=\max_{\chi\in
\widehat{\operatorname{Gal}(L/F)}}
N_{F/\mathbb Q}\mathfrak f_\chi
\]
be the maximum Hecke-character
conductor norm.

The group of units congruent to $1\bmod\ell$ modulo $\ell^n$ has $\ell$-power order dividing $m$. Together with the additive characters and the cyclotomic factor this shows that $\operatorname{Gal}(L/F)$ has exponent dividing $M=m(q-1)$.

At a finite place $v\mid p$, a character
of exponent dividing $M$ is trivial
sufficiently deep in the principal
unit filtration.

Using the $p$-adic logarithm beyond
depth $e_v/(p-1)$ gives the conservative
conductor-exponent bound
\[
\boxed{
a_v(\chi)\le
2+\left\lfloor\frac{e_v}{p-1}\right\rfloor
+e_vv_p(M).
}
\tag{2}
\]

For $p\nmid M$, only tame ramification
can remain.

All ramification of $L/F$ lies above
the fixed rational primes in
$S_0\cup\{q\}$.

Using
\[
\sum_{v\mid p}e_vf_v=d
\]
and $M\le m(q-1)$, equation (2) gives
\[
\boxed{
\log Q\ll_{E,\Gamma,\ell,q,r}
d\log(2m)
\ll_{E,\Gamma,\ell}
m\log(2m).
}
\tag{3}
\]

Importantly, these estimates concern
the degree-$O(m)$ base field $F$ and
the maximum conductor of an abelian
extension over $F$.

They do not substitute the much larger
discriminant of the full division-point
Galois field.

\section{Unconditional degree-one Chebotarev}

We use Thorner--Zaman
\cite[Theorem 1.1]{TZ}.

For a finite abelian extension $L/F$,
their theorem gives, for every prescribed
Frobenius element, a prime ideal
$\mathfrak p$ of $F$ of absolute degree
one, with rational prime norm $p$,
satisfying
\[
\boxed{
p\ll
|D_F|^{694}Q^{521}
+
|D_F|^{232}Q^{367}d^{290d}.
}
\tag{4}
\]

The implied constant is absolute.
No generalized Riemann hypothesis is
assumed.

The degree-one requirement is essential:
it yields a rational prime
$p=N_{F/\mathbb Q}\mathfrak p$,
rather than a prime of higher residue
degree.

Substituting (1) and (3) into (4), and
using
\[
d\log(2d)\ll m\log(2m),
\]
we obtain
\[
\boxed{
\log p\le C_{E,\Gamma,\ell}
m\log(2m).
}
\tag{5}
\]

Apply this result to the $r$ classes
specified in the preceding section.

It produces distinct good rational
primes $p_1,\ldots,p_r$ satisfying (5).

For any $a\ne0\bmod m$, write
$a=\ell^kb$ and choose an index $i$
for which $b_i$ is a unit modulo $\ell$.

Lemma \ref{lem:detection} gives
\[
\Phi_{m,p_i}(a)\ne0.
\]

Thus the joint map has trivial kernel,
proving Theorem \ref{thm:main}.

\section{An exact synchronized-average corollary}

The theorem also supplies a finite-prime
realization of the synchronized
divisibility average.

For a good prime $p$, put
\[
k_p(m)=|E(\mathbb F_p)[m]|
\]
and
\[
N_{p,m}(a)=
\#\{Q\in E(\mathbb F_p):
[m]Q=\sum_i a_i\overline P_i\}.
\]

Then
\[
N_{p,m}(a)=
k_p(m)\,
\mathbf1_{\{\Phi_{m,p}(a)=0\}}.
\]

Let $\mathcal P_X$ be the primes $p\le X$ of good reduction other than $\ell$. Define
\[
\mathcal A_m(X)=
\frac1{m^r}
\sum_{a\in(\mathbb Z/m)^r}
\prod_{p\in\mathcal P_X}
N_{p,m}(a).
\]

The same coefficient vector $a$ is used
at every prime.

Character counting gives the exact
identity
\[
\boxed{
\mathcal A_m(X)=
\frac{\displaystyle\prod_{p\in\mathcal P_X}k_p(m)}
{|\operatorname{im}\Phi_{m,X}|}.
}
\]

Consequently,
\[
R_m(X):=
\frac{\prod_{p\in\mathcal P_X}k_p(m)}
{\mathcal A_m(X)}
=
|\operatorname{im}\Phi_{m,X}|.
\]

For
\[
X\ge\exp(Cm\log(2m)),
\]
Theorem \ref{thm:main} implies
\[
\boxed{R_m(X)=m^r.}
\]

Choose an $\ell$-power $m(X)$ of size
\[
m(X)\asymp
c\frac{\log X}{\log\log X},
\]
with $c>0$ sufficiently small relative
to the constant in Theorem
\ref{thm:main}.

Then $X$ exceeds the detection bound
for all sufficiently large $X$, and
\[
\log R_{m(X)}(X)
=r\log m(X).
\]

Therefore
\begin{corollary}
\[
\boxed{
\lim_{X\to\infty}
\frac{\log R_{m(X)}(X)}
{\log\log X}=r.
}
\]
\end{corollary}

This is a finite-reduction formula for
the rank of the specified arithmetic
subgroup $\Gamma$, under
$(\mathrm K_\ell)$.

Its logarithmic scale is chosen using
the effective Kummer--Chebotarev bound.

\section{Necessity of the hypothesis}

The full first-level Kummer hypothesis
cannot be replaced by mere linear
independence in $\Gamma$.

For example, take a nontorsion
$Q\in E(\mathbb Q)$ and
\[
\Gamma=\langle\ell Q\rangle.
\]

The class of $\ell Q$ is nonzero in
$\Gamma/\ell\Gamma$, but
\[
\overline{\ell Q}\in
\ell E(\mathbb F_p)
\]
at every good prime $p$.

Hence the joint reduction map for
$m=\ell$ is never injective.

The full Kummer hypothesis excludes
this failure. It also excludes more
subtle dependencies arising from the
intersection of point-division fields
with the infinite torsion tower.

The theorem does not assert that
$(\mathrm K_\ell)$ holds for every
elliptic curve, every subgroup, or
every prime $\ell$.

For applications, it must be verified
or deduced from an appropriately
strong Kummer-image theorem.

\section{Relation to published results and scope}

The underlying Kummer-image problem
is classical. Bertrand
\cite{Bertrand} proves full
first-level Kummer images for
sufficiently large primes under
independence hypotheses.

Lombardo--Tronto
\cite{LT} establish effective,
uniform degree estimates for
elliptic Kummer extensions.
Those estimates need not give exact
surjectivity at a specified $\ell$
over the entire infinite torsion field.

Banaszak--Bara\'nczuk
\cite{BB} study reduction maps,
local--global principles and
independence of Kummer towers.

The closest quantitative local--global
comparisons include Sha--Shparlinski
\cite{SS}, on effective detection of
linear dependence, and
Dvornicich--Paladino \cite{DP}, on
effective hypotheses for local--global
divisibility through division fields.

Those works provide substantial
precedent for Kummer--Chebotarev
methods. The present statement isolates
a particular covector-line construction
and its conditional $O(m\log m)$
least-witness exponent; no claim of
publication priority is made.

The critical economy is the use of a
degree-$O(m)$ field fixing a primitive
covector line, rather than requiring
the full torsion field as the base of
the least-prime estimate.

The result is strictly arithmetic.
It neither assumes nor proves
\[
\operatorname{rank}E(\mathbb Q)
=
\operatorname{ord}_{s=1}L(E,s).
\]

In particular, the synchronized rank
limit does not arise from a complex
$L$-function pole or a regulator
identity. Any such comparison would
require additional arithmetic--analytic
input.

\begin{thebibliography}{9}

\bibitem{TZ}
J. Thorner and A. Zaman,
\emph{An explicit bound for the least prime
ideal in the Chebotarev density theorem},
Algebra \& Number Theory \textbf{11}
(2017), no.~5, Theorem 1.1.
\href{https://doi.org/10.2140/ant.2017.11.1135}{Published article}.

\bibitem{Bertrand}
D. Bertrand,
\emph{Kummer theory on the product of an
elliptic curve by the multiplicative group},
Glasgow Mathematical Journal \textbf{22}
(1981), 83--88.
\href{https://doi.org/10.1017/S001708950000450X}{Published article}.

\bibitem{LT}
D. Lombardo and S. Tronto,
\emph{Effective Kummer theory for elliptic
curves},
International Mathematics Research Notices
(2022), no.~22, 17662--17712.
\href{https://doi.org/10.1093/imrn/rnab216}{Published article}.

\bibitem{BB}
G. Banaszak and S. Bara\'nczuk,
\emph{On reduction maps and arithmetic
dynamics of Mordell--Weil type groups},
Acta Arithmetica \textbf{210}
(2023), 155--174.
\href{https://doi.org/10.4064/aa220801-17-5}{Published article}.

\bibitem{SS}
M. Sha and I. E. Shparlinski,
\emph{Effective results on linear dependence
for elliptic curves},
Pacific Journal of Mathematics
\textbf{295} (2018), no.~1.
\href{https://doi.org/10.2140/pjm.2018.295.123}{Published article}.

\bibitem{DP}
R. Dvornicich and L. Paladino,
\emph{On the division fields of an elliptic
curve and an effective bound to the
hypotheses of the local-global divisibility},
International Journal of Number Theory
(2022).
\href{https://doi.org/10.1142/S1793042122500798}{Published article}.

\end{thebibliography}
\end{document}
